\section{Conclusion}
\label{sec:conclusion}

The full run gives different answers for the two research questions. For RQ1, GPT-5.4-mini reached a skeleton cosine median of 0.853 on 261 generated scenarios. This is close to the 0.85 target, but the Wilcoxon test returned $p=0.355$. We fail to reject H0. The generated files often follow the reference structure, but the result is not strong enough to claim that the model clearly passes the semantic threshold. The interpretation is not that the semantic output is poor. A median near 0.85 suggests that the generated skeletons are often close to the expert-written files. The issue is statistical support for a stronger claim: the full sample does not show enough evidence that the model is above the threshold rather than merely near it.

For RQ2, the result is clearer. Every generated feature file passed \texttt{behave --dry-run}. The executable syntax rate is 100\%, and the exact binomial test gives $p < 0.001$, so we reject H0 for executable syntax. In this benchmark, the pinned GPT-5.4-mini setup is reliable at producing parser-ready Gherkin when it receives a same-domain few-shot example. This finding matters because syntax cleanup is a real cost in BDD automation. If a generated feature file fails at the parser stage, the team cannot even start reviewing whether the scenario is useful.

The safest use case is drafting, not automatic acceptance. A QA engineer could use the model output as a first feature-file draft and spend less time fixing parser errors. The reviewer still has to check whether the expected outcome, edge cases, and project-specific business rules are present. This is especially important for domains such as finance, where one changed amount or date can change the meaning of the scenario. Future work should add retrieval from step-definition files and nearby Gherkin examples, compare the same 261 cases against open-source models such as Llama-family and DeepSeek-family models, and add human or mutation-style review for business intent.
